\chapter{无锁数据结构}

先看一个 Treiber stack 事故。线程 T1 读取 \code{head} 指向节点 A，还没来得及 CAS；线程 T2 弹出 A，释放它，又从对象池分配了一个新节点，地址恰好仍是 A；T1 恢复后看到 \code{head} 的位模式仍等于 A，于是 CAS 成功，把一个生命周期已经变化的节点当成旧节点使用。这里没有 mutex 阻塞，也没有普通数据竞争，atomic 和 CAS 都用上了，但程序仍然错了。CAS 比较的是值，不理解对象的历史。

本章的目标不是把“无锁”塑造成更高级的写法，而是把它还原成工程合同：进展保证是什么，线性化点在哪里，内存序发布了什么，失败重试有什么成本，节点什么时候能释放，ABA 如何防护，关闭和等待由谁处理。无锁数据结构的难点通常不在少写一把锁，而在完整生命周期协议。

\begin{learninggoals}
学完本章后，读者应能完成八件事：
\begin{enumerate}
  \item 区分 blocking、obstruction-free、lock-free 和 wait-free，并说明它们不等于低尾延迟。
  \item 为并发结构指出线性化点，包括成功路径、满/空失败路径和关闭路径。
  \item 推导 SPSC ring buffer 的所有权、不变量、release/acquire 发布和测试边界。
  \item 说明 SPSC、MPSC、SPMC、MPMC 是不同问题，不能把单写者算法冒充多写者算法。
  \item 用 Treiber stack 的 ABA 事故解释为什么地址相同不代表对象身份相同。
  \item 比较 tagged pointer、hazard pointer、epoch reclamation、RCU 和 shared ownership 的成本与适用场景。
  \item 设计无锁结构测试矩阵，覆盖唯一 ID、容量边界、随机交错、CAS 失败、节点复用和回收滞留。
  \item 在 Compute Systems Engine 中填写 \code{LockFreeLifecycleReport}，决定何时保留 mutex，何时采用 SPSC，何时使用成熟库。
\end{enumerate}
\end{learninggoals}

\section{先问是否需要无锁}

Blocking 结构可能因为一个线程持锁后暂停而阻塞其他线程。Lock-free 算法保证系统整体持续前进，但单个线程可能反复失败。Wait-free 算法保证每个线程在有限步内完成，通常实现代价更高。Obstruction-free 只在无竞争时保证进展。这些术语不是装饰，它们定义的是竞争和失败时的承诺。

\topic{无锁不是默认优化按钮}
如果临界区短、竞争不高、状态复杂或需要阻塞等待，mutex 加 condition variable 往往更清楚、更省 CPU，也更容易维护。无锁结构适合边界窄、热点明确、生产者消费者模型稳定、容量和生命周期可证明的路径。高竞争 CAS 循环可能烧满 CPU，产生 cache line 迁移和尾延迟；mutex 在严重竞争时让线程睡眠，反而更节能。

工程决策要从 profile 和语义出发。队列操作只占端到端 2% 时，把它改成复杂 MPMC 无锁队列通常不划算。一个 SPSC 管线每秒传递大量小 batch，队列确实是热点，才值得引入无锁 ring。无锁不是荣誉标签，而是带有高测试成本和高维护成本的工具。

\topic{进展保证要写具体}
“没有 mutex”不等于 lock-free。若算法等待某个槽位从 Writing 变成 Ready，而写该槽位的线程崩溃后无人修复，其他线程可能一直等。若 push 每次都调用可能阻塞的通用 allocator，整体也不能简单宣称 lock-free。进展保证要写清依赖条件：是否允许动态分配，是否允许线程取消，是否会在队列满时返回失败，是否把阻塞等待放在结构外层。

\topic{进展保证和尾延迟是两件事}
Lock-free 保证的是系统整体有线程能前进，不保证某个请求低延迟。一个线程可能在 CAS 循环里失败很多次，业务请求尾延迟仍然很差。Wait-free 保证每个线程有限步完成，但常常需要更多元数据、更多内存和更复杂的证明。Blocking mutex 在竞争严重时让线程睡眠，尾延迟未必总比 CAS 自旋差。工程报告不能只写“lock-free 所以快”，必须记录 CAS 失败次数、重试时间、CPU 占用、满/空返回次数和 p99/p999 延迟。

\begin{center}
\begin{tabular}{p{0.20\linewidth}p{0.30\linewidth}p{0.38\linewidth}}
\toprule
进展术语 & 保证 & 工程解释 \\
\midrule
blocking & 持锁线程停住可能阻塞别人 & 容易证明，适合复杂状态和等待 \\
obstruction-free & 无竞争时能完成 & 需要外部退避或仲裁，单独价值有限 \\
lock-free & 总有某个线程能前进 & 单个线程可能饥饿，尾延迟仍需测量 \\
wait-free & 每个线程有限步完成 & 证明和实现成本最高，接口通常更窄 \\
\bottomrule
\end{tabular}
\end{center}

\topic{先定义失败语义}
无锁结构常把阻塞变成失败返回。队列满时 \code{try_push} 返回 false，队列空时 \code{try_pop} 返回空，这不是异常，而是接口语义。调用者必须决定：重试、退避、切换队列、阻塞在外层 Channel、还是向上游施加背压。若上层没有策略，所谓 non-blocking 只会变成忙等烧 CPU。

\begin{lstlisting}[language=C++,caption={try 接口必须配套外层策略}]
bool submit_with_backoff(SpscIntQueue& queue, std::int32_t value) {
  for (std::int32_t attempt = 0; attempt < 8; ++attempt) {
    if (queue.try_push(value)) {
      return true;
    }
    std::this_thread::yield();
  }
  return false;
}
\end{lstlisting}

这段代码不是推荐最终实现，而是提醒读者：\code{try_push} 的 false 是系统设计的一部分。生产系统可能用 semaphore 睡眠、事件循环注册可写、或者把失败转成背压指标。不能在热循环里无限 \code{while (!try_push()) {}}，否则下游慢会变成上游空转。

\section{SPSC Ring：最合适的起点}

单生产者单消费者队列是学习无锁的最佳起点，因为所有权非常清楚：只有生产者写 \code{tail} 和待发布槽位，只有消费者写 \code{head} 和待释放槽位。生产者读取 \code{head} 判断空间，消费者读取 \code{tail} 判断数据。每个索引只有一个写者，这比多个线程 CAS 同一个端点简单得多。

\begin{systemdiagram}{SPSC ring：单写者约束让协议变简单}
\begin{pdfdiagram}
\node[book title label] at (0,0.72) {producer};
\node[book title label] at (3.05,0.72) {ring slots};
\node[book title label] at (6.25,0.72) {consumer};
\node[book title label] at (9.0,0.72) {约束};

\node[book data,text width=2.1cm] (p) at (0,0) {写 slot\\release tail};
\node[book purple node,text width=2.5cm] (ring) at (3.05,0) {fixed buffer\\one empty slot};
\node[book data,text width=2.1cm] (c) at (6.25,0) {acquire tail\\读 slot};
\node[book teal node,text width=2.25cm] (own) at (9.0,0) {tail 只由 P 写\\head 只由 C 写};

\node[book amber node,text width=2.1cm] (h) at (3.05,-1.35) {head\\consumer writes};
\node[book amber node,text width=2.1cm] (t) at (6.25,-1.35) {tail\\producer writes};
\node[book red node,text width=2.55cm] (bad) at (9.0,-1.35) {两个 producer\\算法失效};

\draw[book strong arrow] (p) -- (ring);
\draw[book strong arrow] (ring) -- (c);
\draw[book arrow] (c) -- (h);
\draw[book arrow] (p) -- (t);
\draw[book dashed arrow] (bad) -- (own);

\node[book label,text width=8.8cm] at (4.55,-2.75) {SPSC 的性能来自角色限制，不来自算法能抵抗任意误用。类型名和接口合同必须写清 Single Producer / Single Consumer。};
\end{pdfdiagram}
\diagramnote{本图要证明的命题是：SPSC 简单的根本原因是单写者所有权。}
\end{systemdiagram}

\topic{不变量和线性化点}
常见 SPSC ring 保留一个空槽区分满和空。若 \code{head == tail}，队列空；若 \code{next(tail) == head}，队列满。生产者写入元素后 release 发布新 \code{tail}，这是 \code{try_push} 对消费者可见的线性化点。消费者 acquire 读取 \code{tail} 后才能读槽位，读取元素后 release 发布新 \code{head}，让生产者知道槽位可复用。

\begin{lstlisting}[language=C++,caption={SPSC ring buffer 的核心实现}]
class SpscIntQueue {
public:
  explicit SpscIntQueue(std::uint64_t capacity)
      : buffer_(static_cast<std::size_t>(capacity + 1), 0) {}

  bool try_push(std::int32_t value) {
    const std::uint64_t tail =
        this->tail_.load(std::memory_order_relaxed);
    const std::uint64_t next_tail = this->next(tail);
    const std::uint64_t head =
        this->head_.load(std::memory_order_acquire);
    if (next_tail == head) {
      return false;
    }
    this->buffer_[static_cast<std::size_t>(tail)] = value;
    this->tail_.store(next_tail, std::memory_order_release);
    return true;
  }

  std::optional<std::int32_t> try_pop() {
    const std::uint64_t head =
        this->head_.load(std::memory_order_relaxed);
    const std::uint64_t tail =
        this->tail_.load(std::memory_order_acquire);
    if (head == tail) {
      return std::nullopt;
    }
    const std::int32_t value =
        this->buffer_[static_cast<std::size_t>(head)];
    this->head_.store(this->next(head), std::memory_order_release);
    return value;
  }

private:
  std::uint64_t next(std::uint64_t index) const {
    const std::uint64_t size =
        static_cast<std::uint64_t>(this->buffer_.size());
    return (index + 1) % size;
  }

  std::vector<std::int32_t> buffer_;
  alignas(64) std::atomic<std::uint64_t> head_{0};
  alignas(64) std::atomic<std::uint64_t> tail_{0};
};
\end{lstlisting}

这段代码只适用于一个生产者和一个消费者。两个生产者同时调用 \code{try_push} 会同时写 \code{tail} 和同一槽位；两个消费者同时调用 \code{try_pop} 会同时写 \code{head}。无锁结构的安全来自约束被遵守，不能靠模糊命名让调用者猜。

\topic{关闭放在外层 Channel}
底层 SPSC ring 可以只提供非阻塞 \code{try_push/try_pop}。若系统需要阻塞等待、关闭、drain、取消和超时，更好的分层是外层 \code{Channel} 负责生命周期，底层 ring 负责快速交换。关闭不是一个可以随便塞进 ring 的 bool，它要唤醒等待者、拒绝新写、排空旧数据，并与线程池 shutdown 协同。第 10 章的 \code{QueueContract} 仍然适用。

\topic{测试 SPSC}
SPSC 测试要覆盖容量 1、2、非二次幂、环绕多轮、生产者快消费者慢、消费者快生产者慢、生产者提前结束、消费者 drain。生产者生成单调序号，消费者检查严格递增和不丢不重。还要故意让两个生产者误用，确认文档和 debug 检查能暴露错误。性能报告记录吞吐、空/满返回次数、CPU 占用和 cache line 对齐影响。

\topic{容量、取模和缓存形状}
SPSC ring 可以使用任意容量配合取模，也可以使用二次幂容量配合 bit mask。二次幂版本更快，但接口要明确实际可用容量和内部数组大小；保留一个空槽时，内部大小为 capacity + 1，二次幂 mask 设计则常把 sequence 或可用容量重新定义。教学版优先使用清楚的 \code{next(index)}，性能版才在测量后改成 mask。优化前要先证明不丢、不重、环绕正确。

\topic{false sharing 和索引缓存行}
SPSC 的 \code{head} 和 \code{tail} 会被两个核心频繁读写。若它们落在同一 cache line，生产者写 \code{tail} 会让消费者所在核心的同一条线失效，消费者写 \code{head} 也会反向影响生产者。把两个索引分开对齐能减少 false sharing，但会增加对象大小。报告要测量对齐前后吞吐和 CPU 占用，不要把 \code{alignas(64)} 当成无条件正确。

\begin{lstlisting}[language=C++,caption={索引对齐表达的是缓存所有权边界}]
struct alignas(64) ProducerIndex {
  std::atomic<std::uint64_t> tail{0};
};

struct alignas(64) ConsumerIndex {
  std::atomic<std::uint64_t> head{0};
};
\end{lstlisting}

\section{SPSC、MPSC、MPMC 不能混讲}

队列难度由写者数量决定。SPSC 每端单写者。MPSC 有多个生产者争用 \code{tail}，消费者独占 \code{head}。SPMC 生产者独占 \code{tail}，多个消费者争用 \code{head}。MPMC 两端都有多个写者，通常需要每槽位序号、CAS 预留、提交状态和复杂失败路径。不要把 SPSC 加一个 CAS 就宣传成 MPMC。

\topic{MPSC 的预留和提交}
MPSC 中，多个生产者可以用 CAS 或 \code{fetch_add} 预留槽位，但预留不等于数据已可读。生产者 A 可能预留了较早槽位却还没写完，生产者 B 预留了较后槽位并写完。消费者不能越过 A 读取 B，除非算法允许乱序并写清语义。因此 MPSC 常需要槽位状态：Reserved、Ready、Consumed，或类似序号。tail 表示预留进度，slot state 表示数据发布进度。

\begin{systemdiagram}{MPSC：预留顺序不等于提交顺序}
\begin{pdfdiagram}
\node[book title label] at (0,0.72) {producer A};
\node[book title label] at (2.8,0.72) {producer B};
\node[book title label] at (5.65,0.72) {slots};
\node[book title label] at (8.55,0.72) {consumer};

\node[book amber node,text width=2.15cm] (a) at (0,0) {reserve slot 10\\写入较慢};
\node[book teal node,text width=2.15cm] (b) at (2.8,0) {reserve slot 11\\先写完};
\node[book data,text width=2.35cm] (slots) at (5.65,0) {slot 10 Reserved\\slot 11 Ready};
\node[book red node,text width=2.15cm] (c) at (8.55,0) {不能越过\\slot 10};

\draw[book strong arrow] (a) -- (slots);
\draw[book strong arrow] (b) -- (slots);
\draw[book dashed arrow] (slots) -- (c);

\node[book label,text width=8.7cm] at (4.25,-1.45) {tail 只说明槽位已被预留。消费者能否读取，要看每个槽位自己的提交状态或 sequence。};
\end{pdfdiagram}
\diagramnote{本图要证明的命题是：多生产者队列必须区分预留进度和可读提交进度。}
\end{systemdiagram}

\topic{MPMC 的槽位序号模型}
经典有界 MPMC 队列常让每个槽位维护 sequence。生产者根据 \code{tail} 找槽位，看到 sequence 表示空闲后 CAS 预留，写数据，再 release 更新 sequence 表示可读。消费者根据 \code{head} 找槽位，看到 sequence 表示可读后 CAS 预留读取，读数据，再更新 sequence 表示下一轮空闲。Sequence 用于区分同一数组位置的不同轮次，避免旧 Empty/Ready 被误认。

这个模型很强，但教学项目不应轻率自研生产级 MPMC。关闭、阻塞等待、动态分配、异常、元素析构、CAS 失败风暴和性能退避都会进来。若确有 MPMC 需求，优先评估成熟库；若只是线程池全局队列，mutex bounded queue 或分片队列加 work stealing 往往更稳。

\topic{有界 MPMC 的槽位状态账本}
MPMC 队列至少有三层状态：全局 \code{tail} 用于生产者抢预留，全局 \code{head} 用于消费者抢读取，每个槽位的 \code{sequence} 用于说明该轮次是空、可读还是已消费。全局索引只解决“谁拿到一个位置”，槽位 sequence 才解决“这个位置当前是哪一轮事实”。缺少 sequence，数组位置环绕后，旧的 Ready/Empty 状态会被误认为新一轮状态。

\begin{center}
\begin{tabular}{p{0.24\linewidth}p{0.34\linewidth}p{0.30\linewidth}}
\toprule
状态字段 & 解决的问题 & 典型错误 \\
\midrule
\code{tail} & 多生产者预留哪个槽位 & 预留后未写完就让消费者读 \\
\code{head} & 多消费者竞争哪个槽位 & 两个消费者读同一元素 \\
\code{slot.sequence} & 区分环绕轮次和可读性 & 旧 Ready 被当成新 Ready \\
\code{slot.storage} & 存放元素生命周期 & 构造失败后槽位状态无法恢复 \\
外层 close & 等待、取消、drain、阻塞策略 & 把关闭塞进每个槽位导致协议爆炸 \\
\bottomrule
\end{tabular}
\end{center}

\topic{成熟库也需要外层合同}
成熟 MPMC 队列通常只承诺队列操作，不承诺你的业务关闭语义。库可能支持 non-blocking \code{try_enqueue}，但不知道你的任务是否可取消；可能支持 blocking pop，但不知道 shutdown 时是否 drain；可能要求元素 noexcept move，或者在内部长期持有内存块。引入库前要把库合同翻译成第 10 章的 \code{QueueContract}，否则 benchmark 再好也可能不适合项目。

\section{Treiber 栈：短代码背后的生命周期}

Treiber stack 的 push/pop 看起来短，正好适合作为反面教材：如果只展示 CAS，不讲回收，就是不完整。push 分配节点，把 \code{node->next} 指向旧 \code{head}，CAS \code{head} 为新节点。pop 读取旧 \code{head} 和 \code{next}，CAS \code{head} 为 \code{next}。线性化点通常是 CAS 成功，但 pop 的节点何时能释放，仍然没有答案。

\begin{lstlisting}[language=C++,caption={Treiber 栈的教学骨架，不包含回收协议}]
struct StackNode {
  std::int32_t value = 0;
  StackNode* next = nullptr;
};

class LockFreeStack {
public:
  void push(StackNode* node) {
    StackNode* observed = this->head_.load(std::memory_order_relaxed);
    do {
      node->next = observed;
    } while (!this->head_.compare_exchange_weak(
        observed,
        node,
        std::memory_order_release,
        std::memory_order_relaxed));
  }

private:
  std::atomic<StackNode*> head_{nullptr};
};
\end{lstlisting}

这段代码不能作为完整栈。它没有 pop，更没有说明节点所有权。节点由调用者拥有、栈拥有、对象池拥有，还是回收系统拥有？pop 返回节点后，其他线程是否还可能持有旧指针？这些问题决定结构能否用于生产。

\topic{ABA 时间线}
ABA 可以这样发生：T1 读取 \code{head=A}，并读取 \code{A->next=B}，随后暂停。T2 pop A，又 pop B，释放 A；对象池把地址 A 分给新节点 C，并 push 到栈顶。T1 恢复后看到 \code{head} 的地址仍是 A，CAS 成功，把 \code{head} 改成旧的 B。地址相同，生命周期已经不同，结构被破坏。

\begin{systemdiagram}{ABA：地址相同不代表对象身份相同}
\begin{pdfdiagram}
\node[book title label] at (0,0.72) {T1};
\node[book title label] at (2.8,0.72) {T2};
\node[book title label] at (5.75,0.72) {allocator};
\node[book title label] at (8.7,0.72) {T1 resumes};

\node[book data,text width=2.1cm] (t1a) at (0,0) {load head\\address A};
\node[book red node,text width=2.2cm] (t2a) at (2.8,0) {pop A\\free A};
\node[book amber node,text width=2.2cm] (alloc) at (5.75,0) {reuse address A\\for new C};
\node[book red node,text width=2.25cm] (cas) at (8.7,0) {CAS sees A\\but identity changed};

\node[book label,text width=8.7cm] at (4.35,-1.45) {CAS 只比较位模式。若旧指针仍可能被线程持有，释放和复用地址就会制造 ABA。};

\draw[book strong arrow] (t1a) -- (t2a);
\draw[book strong arrow] (t2a) -- (alloc);
\draw[book strong arrow] (alloc) -- (cas);
\end{pdfdiagram}
\diagramnote{本图要证明的命题是：无锁节点从结构中摘除后，不能立即等同于无人访问。}
\end{systemdiagram}

\topic{Tagged pointer 和 generation}
带标签指针把指针或索引与 generation 一起比较。地址 A 回到 \code{head} 时，generation 已变化，CAS 不会误认为状态没变。教学中更推荐用数组索引加 generation，而不是偷指针低位，因为后者依赖对齐、地址空间和平台 ABI。Generation 思想也适用于对象池、句柄系统和分布式 epoch：位置相同，不代表身份相同。

\begin{lstlisting}[language=C++,caption={标签句柄的语义形状}]
struct TaggedIndex {
  std::uint32_t index = 0;
  std::uint32_t generation = 0;
};
\end{lstlisting}

\topic{对象池句柄比裸指针更适合教学}
教学项目可以用固定数组对象池和 \code{TaggedIndex} 句柄模拟节点身份。数组 index 说明位置，generation 说明该位置当前是哪一代对象。释放槽位时 generation 增加，旧句柄自然失效。这个模型比裸指针更容易测试 ABA，因为读者能打印 index/generation，也能故意把对象池缩小到 2 个槽位快速复用位置。

\begin{lstlisting}[language=C++,caption={对象池句柄检查把身份问题显式化}]
struct PoolHandle {
  std::uint32_t index = 0;
  std::uint32_t generation = 0;
};

struct PoolSlot {
  std::uint32_t generation = 0;
  StackNode node{};
};

bool is_live(const PoolSlot& slot, PoolHandle handle) {
  return slot.generation == handle.generation;
}
\end{lstlisting}

真正生产级无锁栈仍需完整回收协议；tagged handle 只是让“同一位置不同身份”变得可见。若 generation 位数太小，长时间运行后可能溢出；若旧句柄能跨很久保存，溢出后仍可能 ABA。因此报告要写 generation 位数、最大复用速率和是否接受溢出风险。

\section{回收协议：Hazard、Epoch、RCU}

无锁链式结构必须把“从结构中摘除”和“释放内存”分开。摘除说明新的入口路径到不了该节点，不说明旧操作没有持有指针。C++ 没有自动垃圾回收，必须显式证明没有线程仍可能解引用旧节点。

\topic{Hazard pointer：先声明，再解引用}
Hazard pointer 的流程是：线程读取共享指针；把该指针发布到自己的 hazard 槽位；重新读取共享指针确认仍然相同；确认后才能解引用。删除者不能直接释放节点，而是放入 retired 列表；当列表足够大时，扫描所有 hazard 槽位，释放没有被任何线程声明保护的节点。

它的优点是精确，某个节点没人保护就能释放；代价是读路径要发布 hazard，释放路径要扫描所有线程槽位。线程退出时必须清空槽位，否则 retired 节点可能永远不能释放。教学项目可以固定 worker 数简化注册，但文档必须说明假设。

\begin{systemdiagram}{Hazard pointer：保护声明必须早于解引用}
\begin{pdfdiagram}
\node[book title label] at (0,0.72) {reader};
\node[book title label] at (3.0,0.72) {hazard slot};
\node[book title label] at (6.0,0.72) {deleter};
\node[book title label] at (9.0,0.72) {retire scan};

\node[book data,text width=2.2cm] (load) at (0,0) {load shared\\pointer A};
\node[book amber node,text width=2.2cm] (pub) at (3.0,0) {publish hazard\\A};
\node[book data,text width=2.2cm] (check) at (0,-1.25) {recheck head\\still A};
\node[book red node,text width=2.2cm] (retire) at (6.0,0) {remove A\\retire list};
\node[book teal node,text width=2.2cm] (scan) at (9.0,0) {scan hazards\\A protected};

\draw[book strong arrow] (load) -- (pub);
\draw[book arrow] (pub) -- (check);
\draw[book dashed arrow] (retire) -- (scan);
\draw[book dashed arrow] (pub) -- (scan);

\node[book label,text width=8.8cm] at (4.55,-2.45) {删除者只把节点移入 retired 列表。只有扫描不到任何 hazard 槽位保护该节点时，才能真正释放。};
\end{pdfdiagram}
\diagramnote{本图要证明的命题是：无锁读者必须先公开自己可能解引用的节点，删除者才能安全延迟释放。}
\end{systemdiagram}

\topic{Epoch：批量延迟回收}
Epoch 回收把时间分成阶段。线程进入无锁操作时记录当前 epoch，退出时离开。删除节点时记录退休 epoch。只有当所有线程都离开旧 epoch 后，旧 epoch 中退休的节点才能释放。它把单节点精确保护换成批量延迟回收。

Epoch 的读路径通常更轻，但害怕暂停线程。若某个线程进入 epoch 后被抢占、阻塞 I/O、执行未知回调或长时间不退出，旧节点无法释放，retired 列表和内存峰值会增长。Epoch 临界区必须短，不应包含阻塞等待和用户回调。这个约束要写进报告。

\topic{回收协议也要有指标}
回收失败不会立刻表现为错误结果，常常表现为内存慢慢上涨。实验必须记录 retired 节点数量、最大退休年龄、epoch 推进次数、阻塞回收的线程 id、hazard 扫描耗时和实际释放数量。没有这些指标，epoch 泄漏会被误判为“数据结构正常但内存有点高”。

\topic{RCU：读多写少的系统合同}
RCU 可以看成读多写少的发布和回收模型。读者进入 read-side 临界区，读取当前版本；写者复制新版本，修改后发布新指针，然后等待所有旧读者离开 grace period，再释放旧版本。它适合路由表、配置快照、只读索引、分片元数据。代价是写路径复制、旧版本延迟回收和系统级读侧合同。

Compute Systems Engine 不需要自研完整 RCU，但应借鉴思想：driver 发布新 route table 或 config snapshot 时，worker 读取不可变版本；旧版本在没有读者后释放。若更新频繁、对象很大或读者可能阻塞很久，RCU 思想就不一定合适。

\begin{center}
\begin{tabular}{p{0.20\linewidth}p{0.28\linewidth}p{0.28\linewidth}p{0.14\linewidth}}
\toprule
回收方案 & 适合场景 & 主要成本 & 项目建议 \\
\midrule
不释放/统一释放 & 教学、一次性对象、短进程 & 内存增长 & 可作 reference \\
Tagged handle & 对象池、槽位复用、ABA 检测 & 版本溢出、句柄检查 & 推荐理解身份 \\
Hazard pointer & 指针访问复杂，回收要较精确 & hazard 发布和扫描 & 第 12 章实验选做 \\
Epoch & 操作短、线程固定、读侧要轻 & 暂停线程阻塞回收 & 可做小模型 \\
RCU & 读多写少不可变快照 & 写路径复制和 grace period & 借鉴设计 \\
\bottomrule
\end{tabular}
\end{center}

\section{无锁和等待、关闭、异常}

很多系统需要无锁 fast path 和可阻塞 slow path。SPSC ring 的 \code{try_push} 满时返回 false；上层可以用 semaphore、condition variable、eventfd 或运行时任务图等待空间。底层无锁结构负责快速状态转换，上层 \code{Channel} 负责 close、drain、cancel、timeout 和唤醒。把所有语义塞进一个“万能无锁队列”，通常会把证明复杂度推高到不可维护。

\topic{异常和元素类型}
无锁热路径最好不抛异常。若生产者 CAS 预留槽位后，元素构造抛异常，槽位状态如何回滚？消费者如何知道槽位不可读？因此教学代码使用 \code{std::int32_t}、指针或预构造对象，不是偷懒，而是隔离同步机制。泛型无锁容器必须写明元素要求，例如 \code{noexcept} move、外部构造好对象、析构无副作用，或使用复杂槽位生命周期管理。

\topic{动态分配会改变进展保证}
无锁算法内部如果调用可能阻塞的通用 allocator，就不能简单宣称整个操作 lock-free。分配器可能加锁、系统调用或在内存压力下阻塞。固定数组、有界对象池、预分配节点和回收批处理，是无锁结构常见配套。性能报告必须包含分配和回收成本，否则只测了 CAS 的一半。

\topic{退避策略是协议的一部分}
CAS 失败后立刻重试，在低竞争下很快；在高竞争下可能让所有线程持续打同一条 cache line。退避策略包括短暂停顿、\code{yield}、指数退避、切换本地队列、批量提交或回到 blocking slow path。退避不是补丁，而是负载策略。报告要记录退避次数和平均重试轮数，说明它降低了 CAS 风暴还是只是掩盖吞吐问题。

\begin{lstlisting}[language=C++,caption={CAS 失败计数和退避让竞争可观察}]
bool try_transition(std::atomic<std::int32_t>& state,
                    std::int32_t from,
                    std::int32_t to,
                    std::atomic<std::uint64_t>& cas_failures) {
  std::int32_t expected = from;
  if (state.compare_exchange_strong(expected,
                                    to,
                                    std::memory_order_acq_rel,
                                    std::memory_order_acquire)) {
    return true;
  }
  cas_failures.fetch_add(1, std::memory_order_relaxed);
  return false;
}
\end{lstlisting}

\section{工业设计参照}

\begin{casebox}
\textbf{Linux RCU 和 ring buffer。}
Linux RCU 把读侧临界区、发布新版本和 grace period 回收分开，购买极轻读路径；代价是系统级约束和复杂调试。内核 ring buffer 也围绕 per-cpu、提交点和可观测性设计。可迁移原则是：读写角色、提交点和回收时机必须分开写清。

\textbf{Chase-Lev work stealing deque。}
Owner 从 bottom push/pop，stealer 从 top steal；大多数路径单写者，只有最后一个元素需要 CAS 仲裁。它购买的是本地性和低竞争窃取；代价是 top/bottom、扩容、内存序和边界竞争复杂。第 14 章应先用 mutex deque 验证策略，再决定是否需要无锁 deque。

\textbf{Moodycamel、Folly、Boost.Lockfree 等用户态库。}
成熟库通常明确支持 SPSC、MPSC、MPMC、有界/无界、阻塞/非阻塞、元素要求和内存回收策略。引入前要读接口合同，不要只看 benchmark。若库没有项目需要的 close/drain 语义，就要外层封装 \code{QueueContract}。

\textbf{LMAX Disruptor 和序号环。}
Disruptor 系列设计把序号、槽位、发布和消费者进度拆开，适合低延迟流水线。它购买的是预分配、顺序访问和清楚的发布序号；代价是模型更专用，消费者拓扑和等待策略要显式设计。可迁移原则是：高性能 ring 的核心不是“数组很快”，而是每个序号的所有权、发布和消费进度都有账本。

\textbf{\code{io_uring} completion queue。}
\code{io_uring} 的提交队列和完成队列把生产者消费者边界做成共享 ring，但业务仍要处理 request id、buffer 生命周期、取消和迟到 completion。它说明无锁或低锁 ring 只能交付事件，不能替你证明业务提交语义。第 15 章的 I/O 管线会复用这个边界。

\textbf{JVM/Go 等有 GC 的运行时。}
垃圾回收能缓解节点回收和 ABA 的一部分风险，因为对象不会在仍可达时释放；但 CAS 线性化点、状态机、竞争失败和队列语义仍然存在。C++ 教材必须把回收显式讲出，因为运行时不会替你保底。
\end{casebox}

\section{项目落地：LockFreeLifecycleReport}

Compute Systems Engine 在本章不应把所有队列改成无锁。更好的目标是建立 \code{LockFreeLifecycleReport}：对每个候选无锁结构，写明使用场景、生产者消费者数量、容量、线性化点、进展保证、内存序、ABA 防护、回收策略、关闭外置方式、测试矩阵、性能对照和回退方案。

\begin{center}
\begin{tabular}{p{0.24\linewidth}p{0.34\linewidth}p{0.30\linewidth}}
\toprule
报告字段 & 要回答的问题 & 示例 \\
\midrule
模型 & SPSC、MPSC、SPMC、MPMC、stack、deque & reader -> parser 是 SPSC \\
容量 & 有界还是无界，满时如何处理 & 满时返回 false，Channel 背压 \\
线性化点 & 成功和失败操作何时生效 & push 发布 tail \\
进展保证 & blocking、lock-free、wait-free，依赖什么 & SPSC try 操作无阻塞分配 \\
内存序 & 哪些 release/acquire 发布数据 & tail 发布槽位内容 \\
回收 & 固定数组、统一释放、hazard、epoch、RCU & SPSC 固定数组无需节点回收 \\
关闭 & close/drain/cancel 在哪一层实现 & 外层 QueueContract \\
指标 & CAS 失败、满空次数、retired 长度、尾延迟 & per-worker 汇总 \\
回退 & 锁版 reference 是否保留 & mutex bounded queue \\
\bottomrule
\end{tabular}
\end{center}

\topic{推荐项目路线}
第一步，保留 mutex bounded queue 作为 reference 和默认实现。第二步，在一对一管线中实现 SPSC ring，用严格测试证明不丢不重。第三步，若出现多生产者热点，优先尝试分片 SPSC 或每 worker 队列；只有 profile 证明需要时，再评估成熟 MPSC/MPMC。第四步，无锁栈只作为 ABA 和回收教学模型，不进入核心路径。第五步，所有无锁实验必须输出 CAS 失败、满空次数、内存峰值和关闭时间。

\topic{选型矩阵}
项目不应该问“要不要无锁”，而应问“哪个边界需要哪种并发结构”。单个生产者到单个消费者的高频数据流，SPSC ring 是合理候选；多个 worker 提交到一个低频控制队列，mutex queue 更稳；工作窃取 deque 适合 owner 本地 push/pop 很多、steal 较少的任务运行时；分布式消息模拟通常更需要可追踪状态和故障注入，而不是极限无锁。

\begin{center}
\begin{tabular}{p{0.22\linewidth}p{0.28\linewidth}p{0.38\linewidth}}
\toprule
场景 & 首选结构 & 理由 \\
\midrule
reader -> parser 一对一流 & SPSC ring + 外层 Channel & 角色固定，热路径清楚 \\
全局任务提交队列 & mutex bounded queue & 关闭、异常、等待语义更重要 \\
每 worker 本地任务 & deque，先锁版再优化 & 本地性比全局无锁更关键 \\
多生产者日志缓冲 & 分片 buffer 或成熟 MPSC & 避免所有线程争一个 tail \\
配置/路由快照 & atomic shared pointer 或 RCU 思想 & 读多写少，不可变版本更清楚 \\
教学 ABA 模型 & 小对象池 + tagged handle & 容易复现和解释身份变化 \\
\bottomrule
\end{tabular}
\end{center}

\begin{labbox}
\textbf{实验：SPSC ring 与 mutex queue 对照。}
实现 \code{SpscIntQueue} 和锁版 \code{BoundedQueue<int32_t>}。用同一 producer/consumer 输入生成 0 到 N-1 的序号，验证不丢、不重、严格递增。扫描容量 1、2、3、64、非二次幂容量，记录吞吐、满/空返回次数、CPU 占用和尾延迟。再故意用两个 producer 调用 SPSC，说明接口约束被破坏后为什么不能保证正确。
\end{labbox}

\begin{labbox}
\textbf{实验：ABA 与回收可视化。}
实现一个只有 2 到 4 个槽位的小对象池，每个槽位带 generation。构造 Treiber stack 的教学模型，让一个线程暂停在读取 head 后，另一个线程 pop、free、reuse 同一 index，再让第一个线程恢复。报告分别展示裸 index、tagged index、统一释放、hazard pointer 小模型和 epoch 小模型的行为差异。验收重点不是吞吐，而是能用日志解释每个节点的身份变化和释放时机。
\end{labbox}

\topic{测试矩阵}
无锁测试必须比锁版更严。测试容量一、容量二、非二次幂、wrap around、多轮循环、随机 \code{yield}、生产者快、消费者快、长时间压力。Treiber/节点结构测试要使用小对象池快速复用地址，制造 ABA 条件；开启 AddressSanitizer/ThreadSanitizer 时要理解工具限制；记录 retired 列表长度，防止回收静默失控。

\begin{sourcereadingbox}
本章源码阅读以“线性化点和回收”为主线。
\begin{enumerate}
  \item Linux RCU：阅读 RCU 文档或 \filepath{include/linux/rcupdate.h}，说明读侧临界区、grace period 和回收如何分工。
  \item Linux ring buffer：阅读 \filepath{kernel/trace/ring_buffer.c} 的高层设计，关注 per-cpu、提交点和 reader/writer 边界。
  \item Chase-Lev deque 或 oneTBB work stealing：找 owner 路径和 thief 路径，说明最后一个元素竞争如何线性化。
  \item 成熟用户态队列库：填写支持模型、有界/无界、关闭语义、元素要求、回收策略和测试状态。
\end{enumerate}
阅读报告要说明哪些思想适合 Compute Systems Engine，哪些复杂度暂时拒绝。
\end{sourcereadingbox}

\begin{rubricbox}
本章满分 100 分：
\begin{enumerate}
  \item 适用范围，15 分：明确 SPSC/MPSC/MPMC 或栈/deque 模型，不用模糊“FastQueue”命名。
  \item 线性化点，20 分：成功、满、空、关闭或失败路径都有可解释的生效点。
  \item 内存序和协议，15 分：release/acquire 与发布的数据对应，relaxed 只用于安全指标或单写者本地读取。
  \item 生命周期与回收，20 分：ABA、tagged handle、hazard、epoch、RCU 或固定数组策略讲清，不能直接 delete 未证明安全的节点。
  \item 测试矩阵，15 分：唯一 ID、不丢不重、容量边界、随机交错、节点复用、sanitizer 和长期压力覆盖。
  \item 性能证据，10 分：有锁版对照、CAS 失败或满空次数、内存峰值和尾延迟记录。
  \item 工业案例迁移，5 分：至少分析一个成熟实现的约束、机制、能力、代价和项目取舍。
\end{enumerate}
\end{rubricbox}

\section{本章小结}

无锁结构的本质不是“没有锁”，而是用更细的原子协议维护并发对象。SPSC ring 依靠单写者所有权和 release/acquire 发布；MPSC/MPMC 需要预留、提交和槽位序号；Treiber 栈展示 CAS 线性化点，也暴露 ABA 和回收难题；hazard pointer、epoch 和 RCU 都是在回答“旧节点何时安全释放”。Lock-free 保证系统整体进展，不保证每个线程低延迟；无锁代码必须可证明、可测试、可回退。

第三部分到这里形成闭环：第 9 章把线程和 worker 状态变得可见；第 10 章用锁和条件变量保护不变量与等待谓词；第 11 章把窄状态迁移成可审查的原子协议；本章告诉读者，若继续走向无锁，就必须把线性化点和生命周期一起带上。第 14 章的任务运行时与工作窃取，会复用这些边界。

\section{本章习题}

\begin{exercisebox}
\begin{enumerate}
  \item 为 \code{SpscIntQueue} 写不变量表：容量、head、tail、满/空判断、发布顺序和单写者约束。指出 \code{try_push} 和 \code{try_pop} 的线性化点。
  \item 用两个 producer 误用 SPSC ring，说明它为什么不再正确。不要只说“文档禁止”，要指出哪个字段变成多写者。
  \item 画出 Treiber stack 的 ABA 时间线，并分别说明 tagged pointer、hazard pointer 和 epoch 如何缓解不同部分的问题。
  \item 设计一个对象池句柄 \code{TaggedIndex}，包含 index 和 generation。说明 generation 何时增加，旧句柄如何失效。
  \item 比较 hazard pointer 和 epoch：读路径成本、释放延迟、暂停线程影响、线程退出处理、适用场景。
  \item 填写一个 \code{LockFreeLifecycleReport}，决定 Compute Systems Engine 的某条一对一管线是否值得从 mutex queue 替换成 SPSC ring。
\end{enumerate}
\end{exercisebox}
